\paragraph{Meta-2 (Discuss evidence availability and failure cases).}

\begin{itemize}[leftmargin=1.2em, itemsep=1pt, topsep=2pt]
\item \textit{``\dots the revision must discuss what happens when no feasible or useful counterfactuals are present. For example, empty-set behavior, fallback strategies, and trade-offs compared with model-centric or generative counterfactual approaches.''}
\end{itemize}
\smallskip

\jp{The response can be more direct.}

\revision{
Since DcCA retrieves counterfactuals from observed data, evidence availability is a fundamental aspect of the data-centric paradigm. We have revised the paper to explicitly discuss (i) the behavior when no feasible counterfactual exists, (ii) fallback strategies, and (iii) the trade-offs compared with model-centric and generative approaches.

\textbf{What happens when no feasible counterfactual exists?}

When no feasible counterfactual exists in the observed data under the specified constraints, DcCA returns the empty set. This is an intentional and informative outcome rather than a failure of the system. It certifies that no observed instance satisfies both the opposite-outcome predicate and the admissibility constraints under the current specification $\omega$. Unlike model-centric or generative approaches, which always synthesize a candidate, DcCA explicitly distinguishes between ``no feasible evidence exists'' and ``a feasible counterfactual exists,'' providing a verifiable answer grounded in the data.

\textbf{What is the fallback strategy?}

When the feasible set is empty or contains fewer than $k$ counterfactuals, DcCA provides a graded fallback strategy. First, it progressively relaxes user-specified soft constraints (e.g., increasing the intervention budget or relaxing non-essential preferences) while explicitly reporting the minimal relaxation required to restore feasibility. If no satisfactory data-grounded answer can be found, DcCA can optionally invoke a model-centric or generative counterfactual method as an explicitly labeled fallback. This clearly distinguishes retrieved counterfactuals that are supported by observed evidence from synthesized alternatives that are not.

\textbf{What are the trade-offs compared with model-centric approaches?}

The two paradigms are complementary rather than competing. Data-centric counterfactual analysis emphasizes realism, feasibility, provenance, and verifiability because every returned counterfactual corresponds to an observed instance. Its limitation is that it is naturally bounded by the coverage of the available data. In contrast, model-centric and generative methods can always produce a counterfactual, but the generated instance may not correspond to any observed data object or feasible intervention. Thus, data-centric methods prioritize trustworthy, evidence-supported recommendations, whereas model-centric methods prioritize coverage by synthesizing hypothetical alternatives.

To clarify these points, we added a new paragraph, ``Evidence availability and failure behavior,'' in Section~2. The revision explicitly describes the empty-set semantics, partial-answer behavior, graded fallback strategy, and the complementary trade-offs between data-centric and model-centric counterfactual analysis.
}

\noindent\textbf{Response.}
Our formulation already makes the failure case precise: by Eq.~(1) an instance is admissible \emph{``if and only if $\omega(\mathbf{x}_q,\mathbf{x})<\infty$,''} so when $\omega(q)=\varnothing$ no feasible, outcome-changing real instance exists. We argue this is an epistemic \emph{advantage}: DcCA \emph{certifies} infeasibility, whereas generative/model-centric methods~\cite{wachter2017counterfactual,mothilal2020explaining,schleich2021geco} always synthesize a candidate that, as we note in Section~1, \emph{``do[es] not necessarily correspond to real data objects or feasible interventions''}~\cite{delaney2021uncertaintyestimationoutofdistributiondetection,mahajan2019preserving}. Our new validity study makes this coverage effect concrete: in Table~\ref{tab:validity}, DcCA returns feasible, in-distribution counterfactuals whenever a matching instance exists, while for $16\%$ of queries no observed instance shares the query's immutable profile---exactly the empty-set/fallback regime---triggering the graded fallback below.

\noindent\emph{Revision to Section~2 (new paragraph ``Evidence availability and failure behavior''):}
\begin{quote}
\textcolor{blue}{When $\omega(q)=\varnothing$, the query has no feasible, outcome-changing counterfactual in $\mathcal{D}$; DcCA returns the empty set, which certifies infeasibility under $\omega$---a verifiable statement that generative methods cannot make. When $0<|\omega(q)|<k$, DcCA returns a partial set with provenance. To improve coverage, DcCA supports a graded fallback: (i) a constraint-relaxation ladder that relaxes soft constraints or raises the cost budget and reports the minimal relaxation restoring feasibility; and (ii) an explicitly flagged hand-off to a model-centric/generative method as an unverified fallback. The trade-off is thus explicit: data-centric retrieval guarantees realism, feasibility, and verifiability but is bounded by the coverage of $\mathcal{D}$, whereas generative methods always return an answer at the cost of feasibility and verifiability; the two are complementary. Coverage is itself measurable and grows with $\mathcal{D}$.}
\end{quote}
